% page 226 of the facsimile (image p-225 of the scan)
% block 167.6 x 242.0 mm ; left margin 8.0 mm
\pgc{-12.700mm}{0.000mm}{
\l{\cel{3}218}
\l{}
\l{\sou{hip}erbolo. I. (\sou{geom.}) Kurvo qua esas la loko di la punti qui}
\l{\cel{3}distas de du punti fixa konstante. - II. (\sou{retor.})\cel{2}Stilo-fi-}
\l{\cel{3}guro qua konsistas ye exajerar la expreso-termino. - DEFIRS.}
\l{}
\l{hiperboloido. (geom.)\cel{2}Surfaco di la grado duesma, kun centro}
\l{\cel{3}unika, produktita da hiperbolo. - DEFIRS.}
\l{}
\l{hiperboreo. Habitanto di la regioni qui esas extrem-norde. -DEF.}
\l{}
\l{\sou{hiperemio.}\cel{2}(\sou{pat}ol.) Prezenteso di tro multa sango en organo}
\l{\cel{3}o parto de organo.}
\l{}
\l{\sou{hiperestezio}. (\sou{patol.})\cel{2}Hiper-sentiveso.}
\l{}
\l{\sou{.hipermetropo.(patol.}) Qua afektesas da ta anomaleso di refrakto}
\l{\cel{3}pro qua la lumo-radii paralela, vice konvergar a la retino,}
\l{\cel{3}formacas sua foko trans ica kande la binokl-efiko ne interve-}
\l{\cel{3}nas. - DEFS.}
\l{}
\l{\sou{hipertonika}. (\sou{biol.}) Dicesas pri la +solutei injektebla di qui la}
\l{\cel{3}koncentreso-grado molekulala superesas olta di la sero san-}
\l{\cel{3}gala.}
\l{}
\l{\sou{hipertrofiar}. (\sou{trans.}) Produktar hiper-kresko che organo o parto}
\l{\cel{3}de organo, sen altero minboniganta. - DEFIRS.}
\l{}
\l{\sou{hipnotar}. (\sou{netrans.})\cel{2}Submisesar a dormo artificala, produktita da}
\l{\cel{3}la vido longa-dura di objekto briloza. - DEFIRS.}
\l{}
\l{\sou{hipo-}. Prefixo qua signifikas \textquotedbl{}sub-, diminutanta, nesuficanta\textquotedbl{}.}
\l{}
\l{\sou{hipocikloido}.- (\sou{geom.)}\cel{2}Kazo en qua la cirkli movanta di la epi-}
\l{\cel{3}cikloido esas extera. - DEFIRS.}
\l{}
\l{\sou{hipodermo}.- (\sou{anat}.)\cel{2}Tisuo sub-derma.}
\l{}
\l{\sou{hipodromo}.\cel{2}I. (en\cel{1}\sou{la epoki antiqua}). Cirko por la kavalo-konkur-}
\l{\cel{3}si e la charo-konkursi. - II. Agro por la kaval-konkursi.}
\l{\cel{3}- DEFIRS.}
\l{}
\l{.hipofizo. (\sou{anat.)}\cel{2}Organo glanda, an la edro infra di la cerebro}
\l{\cel{3}ed inkluzita en la kavajo di la osto +sfenoida, infre. - DEFIRS}
\l{}
\l{\sou{hipogastro}. (\sou{anat.}) Infra parto di la ventro. - DEFIS.}
\l{}
\l{\sou{hipogeo}. (\sou{arkeol.}) Sepulteyo masonura sub-sula. - DEFI.}
\l{}
\l{\sou{hipogrifo}. (\sou{mitol.}) Animalo fantastika, mi-kavala, mi-grifona,}
\l{\cel{3}rezultinta de la kopulaco da grifono kun kavalino.-DEFIRS.}
\l{}
\l{\sou{hipokampo.}\cel{1}(\sou{zool.)}\cel{2}Genero de fishi, tipo de la familio \textquotedbl{}hipo-}
\l{\cel{3}kampidei\textquotedbl{}, kun formo longa, korpo dividita ad-ringe, e po-}
\l{\cel{3}ligonala-secione, kovrita ye plaki harda. - L.\cel{1}\sou{hippocampus}}
\l{\cel{3}- EFIS.}
}
